\chapter{Die sechsundzwanzig Bände}

\section{Übersicht}

\begin{merksatz}
\boxtitel{Alle Bände mit Gebiet und Umfang}
\begin{center}
{\small
\renewcommand{\arraystretch}{1.32}
\begin{longtable}{@{}P{1.4cm}P{4.3cm}P{4.5cm}r@{}}
\toprule
\textbf{Band} & \textbf{Titel} & \textbf{Gebiet} & \textbf{S.}\\
\midrule
\endhead
\multicolumn{4}{@{}l}{\rule{0pt}{3.4ex}\textbf{\color{mittel}Grundlagen}}\\[2pt]
Band 0 & \textit{Die Logik des Schalters} & Grundlagen des Rechnens & 226\\
Band 1 & \textit{Die Ordnung der Daten} & Algorithmen und Datenstrukturen & 216\\
Band 2 & \textit{Die Grenze der Maschine} & Berechenbarkeit und formale Sprachen & 176\\
Band 3 & \textit{Der Preis der Berechnung} & Komplexitätstheorie & 166\\
\multicolumn{4}{@{}l}{\rule{0pt}{3.4ex}\textbf{\color{mittel}Die Maschine und ihre Schichten}}\\[2pt]
Band 4 & \textit{Die Anatomie der Geschwindigkeit} & Rechnerarchitektur & 166\\
Band 5 & \textit{Der Wettlauf der Prozesse} & Betriebssysteme und Nebenläufigkeit & 186\\
Band 6 & \textit{Die Verwandlung des Textes} & Programmiersprachen und Übersetzerbau & 188\\
\multicolumn{4}{@{}l}{\rule{0pt}{3.4ex}\textbf{\color{mittel}Information, Sicherheit, Verbindung}}\\[2pt]
Band 7 & \textit{Das Maß der Nachricht} & Informationstheorie & 158\\
Band 8 & \textit{Die Mathematik des Misstrauens} & Kryptographie & 174\\
Band 9 & \textit{Die Einigung der Getrennten} & Netzwerke und verteilte Systeme & 166\\
Band 10 & \textit{Das Gedächtnis der Systeme} & Datenbanken und Information Retrieval & 162\\
Band 11 & \textit{Das Muster der Verbindungen} & Graphen und Netzwerkanalyse & 164\\
\multicolumn{4}{@{}l}{\rule{0pt}{3.4ex}\textbf{\color{mittel}Schließen und Lernen}}\\[2pt]
Band 12 & \textit{Der Raum der Möglichkeiten} & Symbolische künstliche Intelligenz & 164\\
Band 13 & \textit{Der Schluss aus dem Beispiel} & Maschinelles Lernen & 170\\
Band 14 & \textit{Das Rechnen des Gehirns} & Computational Neuroscience & 174\\
Band 15 & \textit{Die Auslöschung der Irrwege} & Quanteninformatik & 156\\
\multicolumn{4}{@{}l}{\rule{0pt}{3.4ex}\textbf{\color{mittel}Rahmen}}\\[2pt]
Begleit\-band & \textit{Der seltsame Kreis} & Selbstbezug, Universalität und Grenze & 124\\
\multicolumn{4}{@{}l}{\rule{0pt}{3.4ex}\textbf{\color{mittel}Herkunft und Praxis}}\\[2pt]
Band 16 & \textit{Die Maschine, die rechnen lernte} & Eine Geschichte des maschinellen Rechnens & 146\\
Band 17 & \textit{Das Werk vieler Hände} & Softwaretechnik & 150\\
Band 18 & \textit{Die Übersetzung der Absicht} & Mensch-Maschine-Interaktion & 142\\
Band 19 & \textit{Der schwächste Punkt} & Sicherheit & 150\\
\multicolumn{4}{@{}l}{\rule{0pt}{3.4ex}\textbf{\color{mittel}Programmieren}}\\[2pt]
Band 20 & \textit{Der zweite Entwurf} & Programmieren in Python & 139\\
Band 21 & \textit{Was darunter liegt} & Programmieren in C & 132\\
Band 22 & \textit{Die Kunst des Weglassens} & Programmieren in Go & 118\\
Band 23 & \textit{Vom Programm zum Werk} & Programme, die zusammenwirken & 122\\
Band 24 & \textit{Die Kunst} & Zwanzig Fälle und ihre Begründung & 108\\
\midrule
\textbf{Summe} & \multicolumn{2}{@{}l}{\textbf{26 Bände}} & \textbf{4143}\\
\bottomrule
\end{longtable}}
\end{center}
\end{merksatz}

\begin{center}
\begin{minipage}{0.88\textwidth}
\kursiv{Die Reihe war auf siebzehn Bände angelegt und ist auf sechsundzwanzig gewachsen. Die neun zusätzlichen --- Geschichte, Softwaretechnik, Mensch-Maschine-Interaktion, Sicherheit und die fünf Programmierbände --- behandeln kein neues Gebiet der Theorie, sondern die Frage, was man mit ihr tut.}
\end{minipage}
\end{center}

\section{Was in welchem Band steht}

Jeder Band ist für sich lesbar. Die folgenden Kurzbeschreibungen sagen, was man vorfindet und wofür man ihn aufschlägt.

\begin{merksatz}
\boxtitel{Grundlagen}
\begin{center}
{\small
\renewcommand{\arraystretch}{1.42}
\begin{tabular}{@{}P{2.2cm}P{6.6cm}@{}}
\toprule
\textbf{Band 0} & Vom Schalter über Gatter und Flipflop bis zur universellen Maschine. Der einzige Band, der wirklich bei null anfängt.\\
\textbf{Band 1} & Sortieren, Suchen, Bäume, Streuspeicher --- mit vollständigen Aufwandsrechnungen statt bloßer Ergebnisse.\\
\textbf{Band 2} & Automaten, Grammatiken, Turingmaschine, Halteproblem. Was grundsätzlich nicht berechenbar ist.\\
\textbf{Band 3} & P, NP, Reduktionen, Approximation. Was berechenbar, aber nicht bezahlbar ist.\\
\bottomrule
\end{tabular}}
\end{center}
\end{merksatz}

\begin{merksatz}
\boxtitel{Die Maschine und ihre Schichten}
\begin{center}
{\small
\renewcommand{\arraystretch}{1.42}
\begin{tabular}{@{}P{2.2cm}P{6.6cm}@{}}
\toprule
\textbf{Band 4} & Fließband, Sprungvorhersage, Zwischenspeicher, Speicherhierarchie --- warum ein Programm schnell oder langsam läuft.\\
\textbf{Band 5} & Prozesse, Fäden, Sperren, Verklemmungen, virtueller Speicher. Enthält das Kapitel zu Wettläufen, auf das Band 22 zurückgreift.\\
\textbf{Band 6} & Von der Zeichenkette zum lauffähigen Programm: Zerteilen, Syntaxbaum, Zwischencode, Optimierung.\\
\bottomrule
\end{tabular}}
\end{center}
\end{merksatz}

\begin{merksatz}
\boxtitel{Information, Sicherheit, Verbindung}
\begin{center}
{\small
\renewcommand{\arraystretch}{1.42}
\begin{tabular}{@{}P{2.2cm}P{6.6cm}@{}}
\toprule
\textbf{Band 7} & Entropie, Kodierung, Kanalkapazität. Die Antwort auf die Frage, wie viel man überhaupt übertragen kann.\\
\textbf{Band 8} & Die Mathematik hinter Schlüsseln und Signaturen --- und die ehrliche Auskunft, woran Systeme stattdessen scheitern.\\
\textbf{Band 9} & Schichten, Protokolle, Uhren, Konsens. Was zwei Rechner voneinander wissen können und was nicht.\\
\textbf{Band 10} & Relationen, Anfrageplanung, Transaktionen, Volltextsuche. Wie Systeme sich erinnern.\\
\textbf{Band 11} & Wege, Zentralität, Gemeinschaften --- und die Prüfung verbreiteter Behauptungen über Netzwerke.\\
\bottomrule
\end{tabular}}
\end{center}
\end{merksatz}

\begin{merksatz}
\boxtitel{Schließen und Lernen}
\begin{center}
{\small
\renewcommand{\arraystretch}{1.42}
\begin{tabular}{@{}P{2.2cm}P{6.6cm}@{}}
\toprule
\textbf{Band 12} & Suche, Logik, Planung, Wissensrepräsentation --- die künstliche Intelligenz vor dem Lernen.\\
\textbf{Band 13} & Regression, Bäume, Netze, Überanpassung. Mit dem Kapitel darüber, was ein Modell nicht kann.\\
\textbf{Band 14} & Neuron, Netz, Kodierung, Konnektom. Wo die Analogie zwischen Gehirn und Rechner trägt und wo nicht.\\
\textbf{Band 15} & Qubit, Verschränkung, Algorithmen, Fehlerkorrektur --- ohne die üblichen Übertreibungen.\\
\bottomrule
\end{tabular}}
\end{center}
\end{merksatz}

\begin{merksatz}
\boxtitel{Rahmen}
\begin{center}
{\small
\renewcommand{\arraystretch}{1.42}
\begin{tabular}{@{}P{2.4cm}P{6.4cm}@{}}
\toprule
\textbf{Begleitband} & Selbstbezug, Universalität und Grenze --- der Faden, der Band 0 mit Band 15 verbindet. Kein Fachgebiet, sondern die Frage, was Systeme über sich selbst wissen können.\\
\bottomrule
\end{tabular}}
\end{center}
\end{merksatz}

\begin{merksatz}
\boxtitel{Herkunft und Praxis}
\begin{center}
{\small
\renewcommand{\arraystretch}{1.42}
\begin{tabular}{@{}P{2.2cm}P{6.6cm}@{}}
\toprule
\textbf{Band 16} & Von der Rechenhilfe zum Rechenzentrum in der Wolke. Warum die Dinge so geworden sind, wie sie sind.\\
\textbf{Band 17} & Entwurf, Prüfung, Durchsicht, Ablösung --- und die Zahlen dazu. Der Band, aus dem die Regeln stammen, die Band 24 auf die Probe stellt.\\
\textbf{Band 18} & Die Kluft zwischen Absicht und System: Bedienbarkeit, Fehlbedienung, Rückmeldung.\\
\textbf{Band 19} & Warum Systeme brechen und woran. Die Vertrauensgrenze aus Kapitel~4 ist der meistzitierte Begriff der ganzen Reihe.\\
\bottomrule
\end{tabular}}
\end{center}
\end{merksatz}

\begin{merksatz}
\boxtitel{Programmieren}
\begin{center}
{\small
\renewcommand{\arraystretch}{1.42}
\begin{tabular}{@{}P{2.2cm}P{6.6cm}@{}}
\toprule
\textbf{Band 20} & Python. Vom Lauffähigen zum Guten --- der zweite Entwurf als Arbeitsform, nicht als Reue.\\
\textbf{Band 21} & C. Was unter den Abstraktionen liegt: Zeiger, Speicher, undefiniertes Verhalten.\\
\textbf{Band 22} & Go. Nebenläufigkeit, ausgeschriebene Fehlerpfade und die Kunst des Weglassens.\\
\textbf{Band 23} & Was geschieht, wenn Programme in mehreren Sprachen zusammenwirken: Zeit, Text, Zahlen, Ausfall an den Nähten.\\
\textbf{Band 24} & Zwanzig Fälle aus einem gewachsenen System, jeder mit Urteil und Begründung --- vierzehn Änderungen, sechs Begründungen für Zurückhaltung.\\
\bottomrule
\end{tabular}}
\end{center}

\begin{center}
\begin{minipage}{0.88\textwidth}
\kursiv{Die fünf Programmierbände sind aufeinander bezogen und behandeln denselben klinischen Anwendungsfall --- ein System aus Messerfassung, Übertragung und Auswertung. Wer sie in der Reihenfolge liest, sieht dasselbe Vorhaben viermal aus verschiedener Höhe.}
\end{minipage}
\end{center}
\end{merksatz}
